% Hình học của gradient clipping: g = (3, 4) có chuẩn 5 bị co về hình tròn bán kính c = 1, giữ hướng.
\begin{tikzpicture}[scale=0.95]
  \draw[gray!30] (-1.5,-1.5) grid (4.5,4.5);
  \draw[-{Stealth}, gray] (-1.5,0) -- (4.7,0) node[right, font=\footnotesize] {$g_1$};
  \draw[-{Stealth}, gray] (0,-1.5) -- (0,4.7) node[above, font=\footnotesize] {$g_2$};
  \fill[lightGreen, opacity=0.7] (0,0) circle (1);
  \draw[mThree, thick] (0,0) circle (1);
  \node[mThree!70!black, font=\footnotesize, anchor=north west] at (0.72,-0.72) {$\lVert g \rVert \le c = 1$};
  \draw[-{Stealth[length=3mm]}, very thick, accentRed] (0,0) -- (\SL{theory/clip/g2/0Raw},\SL{theory/clip/g2/1Raw})
    node[above right, font=\small] {$g$, $\lVert g\rVert = \SL{theory/clip/g2_norm}$};
  \draw[-{Stealth[length=3mm]}, very thick, ptitblue] (0,0) -- (\SL{theory/clip/g2_clipped/0Raw},\SL{theory/clip/g2_clipped/1Raw})
    node[left=1mm, font=\small, fill=white, inner sep=1pt] {$g \cdot \dfrac{c}{\lVert g\rVert}$};
  \draw[dashed, gray] (\SL{theory/clip/g2_clipped/0Raw},\SL{theory/clip/g2_clipped/1Raw}) -- (\SL{theory/clip/g2/0Raw},\SL{theory/clip/g2/1Raw});
  \node[font=\footnotesize, align=left, anchor=west] at (1.6,1.2) {cùng hướng,\\chỉ đổi độ dài};
\end{tikzpicture}
